\chapter*{Remerciements}
\addcontentsline{toc}{chapter}{Remerciements}
\markboth{REMERCIEMENTS}{}

Au terme de ce projet de fin d'etudes, je tiens a exprimer ma profonde gratitude
a toutes les personnes qui ont contribue, de pres ou de loin, a la reussite de
ce travail.

Je tiens, en premier lieu, a remercier vivement mon encadrant universitaire,
\textbf{Mme/M. \dots\dots\dots\dots} (\textit{nom à compléter}), pour la qualite de son
encadrement, sa disponibilite, ses precieux conseils et la confiance qu'il
m'a accordee tout au long de la realisation de ce projet. Sa rigueur et ses
orientations ont enormement contribue a faconner non seulement le contenu de
ce rapport, mais aussi la demarche avec laquelle j'ai aborde chacune de ses
etapes.

J'adresse egalement mes sinceres remerciements a mon encadrant professionnel,
\textbf{Mme/M. \dots\dots\dots\dots} (\textit{nom à compléter}), pour le temps qu'il m'a consacre,
pour la richesse des echanges que nous avons eus, et pour son ouverture aux
besoins reels du marche tunisien des petites et moyennes entreprises.

Je remercie tres respectueusement les membres du jury qui ont accepte
d'evaluer ce travail, ainsi que l'ensemble du corps professoral de
\textbf{l'Institut Superieur des Mathematiques Appliquees et d'Informatique
de Kairouan} pour la qualite de la formation dispensee tout au long de mon
cursus.

Je tiens a remercier tout particulierement les proprietaires de cliniques,
de laboratoires et d'autres petites entreprises tunisiennes qui ont accepte
de tester Nova~AI sur leurs propres clients pendant les sprints de
developpement. Leurs retours, leurs critiques et leurs suggestions ont
oriente le projet bien plus surement que n'aurait pu le faire la seule
theorie academique.

Enfin, je remercie ma famille et mes amis pour leur soutien indefectible
dur